\chapter{Fuzzy Sublocales, Endpoint Quotients, and Open--Closed Gluing}
\label{ch:fuzzy-sublocales-open-closed-gluing}

The preceding chapters developed categorical fuzzy frames, fuzzy locales, and
the fuzzy Sierpiński classifier.  We now use that structure to develop
sublocales, endpoint quotients of fuzzy opens, generated open and
closed-complement sublocales, attachment functors, and Artin gluing.

There is one important structural distinction throughout this chapter.  The
general scalar-compatible fuzzy theory is not the same as the ordinary
underlying-frame theory.  In an ordinary frame, an element \(a\in L\) generates
the two familiar nuclei
\[
    o_a(x)=a\Rightarrow x,
    \qquad
    c_a(x)=a\vee x.
\]
These formulas remain correct in the ordinary, idempotent, and crisp
specializations.  They are not, however, the general scalar-compatible fuzzy
formulas.  For non-idempotent bases such as Łukasiewicz or product truth values,
the map \(a\Rightarrow-\) need not be compatible with scalar action and hence
need not define a fuzzy nucleus.

Accordingly, the general construction is expressed invariantly.  A fuzzy open
\[
    U\hookrightarrow 1_{\mathcal A}
\]
of a categorical fuzzy frame \(\mathcal A\) is classified by a fuzzy-frame
morphism
\[
    \chi_U^\ast:\Omega_{\mathbb V}\to\mathcal A
\]
from the fuzzy Sierpiński frame.  The generated true and false endpoint
quotients are obtained by base change along the two endpoint quotients of
\(\Omega_{\mathbb V}\):
\[
    \mathcal A_U
    =
    \mathcal A\amalg_{\Omega_{\mathbb V}}\Omega_{\mathbb V}^{\top},
    \qquad
    \mathcal A_Z
    =
    \mathcal A\amalg_{\Omega_{\mathbb V}}\Omega_{\mathbb V}^{\bot}.
\]
In the posetal scalar-compatible reflection these are described by generated
fuzzy nuclei.  The false endpoint has the explicit formula
\[
    x\longmapsto a\vee x.
\]
The true endpoint is the least scalar-compatible fuzzy nucleus forcing \(a\) to
be top; in general it has no pointwise formula \(a\Rightarrow-\).

The general fuzzy theory therefore provides a canonical gluing comparison
\[
    \mathcal A\longrightarrow \operatorname{Gl}(\Theta_U),
\]
where
\[
    \Theta_U=i_U^\ast(j_U)_\ast:
    \mathcal A_U\to\mathcal A_Z.
\]
This comparison is not an equivalence in general.  It becomes an equivalence in
the ordinary/idempotent/crisp specialization, where the classical
open--closed Artin gluing theorem is recovered.

\section{Fuzzy sublocales as left-exact reflective quotients}
\label{sec:fuzzy-sublocales-as-lex-reflective-quotients}

Let \(X\) be a fuzzy locale, and write
\[
    \mathcal O(X)=\mathcal A
\]
for its categorical fuzzy frame.

\begin{definition}[Fuzzy sublocale]
\label{def:fuzzy-sublocale}
A \emph{fuzzy sublocale} of \(X\) is presented by a left-exact reflective
quotient of categorical fuzzy frames
\[
    q\dashv r:
    \mathcal A
    \rightleftarrows
    \mathcal B,
\]
where \(r:\mathcal B\to\mathcal A\) is fully faithful and \(q\) is a morphism
of categorical fuzzy frames.

The fuzzy frame of opens of the presented sublocale is
\[
    \mathcal O(S):=\mathcal B.
\]
The associated idempotent localization modality on \(\mathcal A\) is
\[
    J:=rq:\mathcal A\to\mathcal A.
\]
\end{definition}

Thus a sublocale is a quotient in the frame direction and an inclusion in the
localic direction.  If
\[
    \mathcal O(X)=\mathcal A
    \qquad\text{and}\qquad
    \mathcal O(S)=\mathcal B,
\]
then the inclusion
\[
    S\hookrightarrow X
\]
is represented contravariantly by
\[
    q:\mathcal A\to\mathcal B.
\]

\begin{proposition}[Reflective quotients are fuzzy frames]
\label{prop:reflective-quotients-are-fuzzy-frames}
Let
\[
    q\dashv r:
    \mathcal A
    \rightleftarrows
    \mathcal B
\]
be a left-exact reflective quotient of a categorical fuzzy frame
\(\mathcal A\).  Then \(\mathcal B\) is again a categorical fuzzy frame.
\end{proposition}

\begin{proof}
Choose a presentation of \(\mathcal A\) as a categorical fuzzy frame
\[
    a\dashv i:
    P_{\mathbb V}(X)
    \rightleftarrows
    \mathcal A,
\]
where \(i\) is fully faithful and \(a\) preserves the specified finite weighted
limits.  Composing with the reflective quotient gives
\[
    qa:
    P_{\mathbb V}(X)\to\mathcal B.
\]
Its right adjoint is
\[
    ir:\mathcal B\to P_{\mathbb V}(X).
\]
Since \(i\) and \(r\) are fully faithful, the composite \(ir\) is fully
faithful.  Since \(a\) and \(q\) preserve the specified finite weighted limits,
so does \(qa\).  Hence
\[
    qa\dashv ir:
    P_{\mathbb V}(X)
    \rightleftarrows
    \mathcal B
\]
is a categorical fuzzy-frame presentation of \(\mathcal B\).
\end{proof}

\begin{definition}[Sublocale order]
\label{def:fuzzy-sublocale-order}
Let \(S\) and \(T\) be fuzzy sublocales of \(X\), presented by quotients
\[
    q_S:\mathcal A\to\mathcal B_S,
    \qquad
    q_T:\mathcal A\to\mathcal B_T.
\]
We write
\[
    S\leq T
\]
if \(q_S\) factors through \(q_T\) in the fuzzy-frame direction:
\[
    q_S\cong \bar q\, q_T
\]
for some fuzzy-frame morphism
\[
    \bar q:\mathcal B_T\to\mathcal B_S.
\]
Equivalently, the localic inclusion of \(S\) factors through the localic
inclusion of \(T\).
\end{definition}

\section{The posetal reflection: scalar-compatible fuzzy nuclei}
\label{sec:posetal-reflection-scalar-compatible-nuclei}

Let
\[
    \Lambda=(L,\leq,\meet,\join,\tensor,\multimap,\bot,\top)
\]
be a fuzzy truth-value object, and let \(A\) be a fuzzy frame in the posetal
reflection.

We write the scalar action of \(\lambda\in\Lambda\) on \(x\in A\) as
\[
    \lambda\tensor x.
\]

\begin{definition}[Fuzzy nucleus]
\label{def:fuzzy-nucleus}
A \emph{fuzzy nucleus} on \(A\) is a monotone map
\[
    j:A\to A
\]
such that:

\begin{enumerate}
    \item \(j\) is inflationary:
    \[
        x\leq jx;
    \]

    \item \(j\) is idempotent:
    \[
        jjx=jx;
    \]

    \item \(j\) preserves the terminal object and binary meets:
    \[
        j(1)=1,
        \qquad
        j(x\wedge y)=jx\wedge jy;
    \]

    \item \(j\) is compatible with scalar action:
    \[
        j(\lambda\tensor jx)
        =
        j(\lambda\tensor x)
    \]
    for all \(\lambda\in\Lambda\) and \(x\in A\).
\end{enumerate}
\end{definition}

\begin{proposition}[Fuzzy nuclei and reflective quotients]
\label{prop:fuzzy-nuclei-reflective-quotients}
In the posetal reflection, fuzzy sublocales of a fuzzy locale with frame \(A\)
are presented by fuzzy nuclei
\[
    j:A\to A.
\]
The corresponding fuzzy frame of opens is the fixed-point fuzzy frame
\[
    A_j=\{x\in A\mid jx=x\}.
\]
The quotient map is
\[
    q_j:A\to A_j,
    \qquad
    q_j(x)=jx,
\]
and the fully faithful right adjoint is the inclusion
\[
    i_j:A_j\hookrightarrow A.
\]
The scalar action on \(A_j\) is
\[
    \lambda\tensor_j x
    :=
    j(\lambda\tensor x).
\]
\end{proposition}

\begin{proof}
Let \(j\) be a fuzzy nucleus.  Since \(j\) is inflationary and idempotent, the
fixed points form a reflective quotient of \(A\), with reflector \(q_j(x)=jx\)
and fully faithful right adjoint the inclusion \(i_j\).

Because \(j\) preserves the terminal object and binary meets, the quotient is
left exact.  The scalar-compatibility equation
\[
    j(\lambda\tensor jx)=j(\lambda\tensor x)
\]
is exactly the condition that the displayed scalar action on fixed points is
well defined and that \(q_j\) preserves scalar action.

Conversely, any left-exact reflective quotient
\[
    q\dashv i:A\rightleftarrows B
\]
with \(i\) fully faithful determines the idempotent
\[
    j=iq:A\to A.
\]
Left exactness and scalar compatibility of the quotient give precisely the
conditions in Definition~\ref{def:fuzzy-nucleus}.
\end{proof}

\begin{proposition}[Sublocale order and nuclei]
\label{prop:sublocale-order-nuclei}
For fuzzy nuclei \(j,k\) on \(A\),
\[
    X_j\leq X_k
    \quad\Longleftrightarrow\quad
    k\leq j
\]
pointwise.
\end{proposition}

\begin{proof}
The sublocale order is factorization order on quotients.  The quotient
\[
    q_j:A\to A_j
\]
factors through
\[
    q_k:A\to A_k
\]
if and only if every \(k\)-closed element is sent consistently to a \(j\)-closed
element.  Equivalently, the localization \(j\) is at least as strong as \(k\),
that is,
\[
    kx\leq jx
\]
for all \(x\in A\).  Hence
\[
    X_j\leq X_k
    \quad\Longleftrightarrow\quad
    k\leq j.
\]
\end{proof}

\begin{lemma}[Meets of fuzzy nuclei]
\label{lem:meets-of-fuzzy-nuclei}
Arbitrary pointwise meets of fuzzy nuclei are fuzzy nuclei.
\end{lemma}

\begin{proof}
Let \((j_\alpha)_\alpha\) be a family of fuzzy nuclei, and define
\[
    jx:=\bigwedge_\alpha j_\alpha x.
\]
Inflationarity and monotonicity are immediate.

For idempotence, first \(x\leq jx\) gives \(jx\leq jjx\).  Conversely, for each
\(\alpha\),
\[
    jx\leq j_\alpha x.
\]
By monotonicity,
\[
    j_\alpha(jx)\leq j_\alpha(j_\alpha x)=j_\alpha x.
\]
Taking the meet over \(\alpha\) gives
\[
    jjx\leq jx.
\]

Preservation of finite meets follows from complete distributivity of the
underlying fuzzy frame:
\[
\begin{aligned}
    j(x\wedge y)
    &=
    \bigwedge_\alpha j_\alpha(x\wedge y)        \\
    &=
    \bigwedge_\alpha (j_\alpha x\wedge j_\alpha y)\\
    &=
    \left(\bigwedge_\alpha j_\alpha x\right)
    \wedge
    \left(\bigwedge_\alpha j_\alpha y\right)    \\
    &=
    jx\wedge jy.
\end{aligned}
\]
The terminal object is preserved because each \(j_\alpha\) preserves it.

For scalar compatibility,
\[
\begin{aligned}
    j(\lambda\tensor jx)
    &=
    \bigwedge_\alpha
    j_\alpha\!\left(\lambda\tensor\bigwedge_\beta j_\beta x\right).
\end{aligned}
\]
Since
\[
    x\leq jx\leq j_\alpha x,
\]
we have
\[
    \lambda\tensor x
    \leq
    \lambda\tensor jx
    \leq
    \lambda\tensor j_\alpha x.
\]
Applying \(j_\alpha\) gives
\[
    j_\alpha(\lambda\tensor x)
    \leq
    j_\alpha(\lambda\tensor jx)
    \leq
    j_\alpha(\lambda\tensor j_\alpha x)
    =
    j_\alpha(\lambda\tensor x).
\]
Hence
\[
    j_\alpha(\lambda\tensor jx)
    =
    j_\alpha(\lambda\tensor x)
\]
for every \(\alpha\), and taking meets yields
\[
    j(\lambda\tensor jx)=j(\lambda\tensor x).
\]
Thus \(j\) is a fuzzy nucleus.
\end{proof}

\section{Endpoint quotients of the fuzzy Sierpiński frame}
\label{sec:endpoint-quotients-fuzzy-sierpinski}

The fuzzy Sierpiński frame
\[
    \Omega_{\mathbb V}
\]
classifies fuzzy opens.  It carries a generic fuzzy open
\[
    \sigma_{\mathbb V}\hookrightarrow 1_{\Omega_{\mathbb V}}.
\]
For every categorical fuzzy frame \(\mathcal A\), fuzzy opens
\[
    U\hookrightarrow1_{\mathcal A}
\]
are classified by fuzzy-frame morphisms
\[
    \chi_U^\ast:\Omega_{\mathbb V}\to\mathcal A.
\]

\begin{definition}[True endpoint quotient]
\label{def:true-endpoint-quotient}
The \emph{true endpoint quotient}
\[
    e_{\top}^{\ast}:
    \Omega_{\mathbb V}\to\Omega_{\mathbb V}^{\top}
\]
is the universal left-exact fuzzy-frame quotient in which the generic open
becomes terminal:
\[
    e_{\top}^{\ast}(\sigma_{\mathbb V})
    \cong
    1_{\Omega_{\mathbb V}^{\top}}.
\]
\end{definition}

\begin{definition}[False endpoint quotient]
\label{def:false-endpoint-quotient}
The \emph{false endpoint quotient}
\[
    e_{\bot}^{\ast}:
    \Omega_{\mathbb V}\to\Omega_{\mathbb V}^{\bot}
\]
is the universal left-exact fuzzy-frame quotient in which the generic open
becomes initial in the quotient.
\end{definition}

\begin{proposition}[Existence of endpoint quotients]
\label{prop:endpoint-quotients-exist}
The endpoint quotients
\[
    \Omega_{\mathbb V}^{\top},
    \qquad
    \Omega_{\mathbb V}^{\bot}
\]
exist as categorical fuzzy frames.
\end{proposition}

\begin{proof}
The fuzzy Sierpiński frame is presented by an enriched
finite-limit-and-scalar sketch with one generic subterminal open.  The true
endpoint quotient is obtained by adjoining the equation identifying this
generic open with the terminal object.

The false endpoint quotient is obtained by adjoining the equation that the
generic open becomes the initial object in the quotient.  In the posetal
reflection this condition is expressed by equality with the quotient bottom,
not by an inflationary nucleus sending the original element literally below
itself.

Adding either endpoint equation gives another enriched
finite-limit-and-scalar sketch localization.  Therefore the resulting
quotients are categorical fuzzy frames.
\end{proof}

\section{Generated true and false endpoint sublocales}
\label{sec:generated-endpoint-sublocales}

Let \(\mathcal A\) be a categorical fuzzy frame, and let
\[
    U\hookrightarrow1_{\mathcal A}
\]
be a fuzzy open classified by
\[
    \chi_U^\ast:\Omega_{\mathbb V}\to\mathcal A.
\]

\begin{definition}[Generated true endpoint fuzzy frame]
\label{def:generated-true-endpoint-frame}
The \emph{generated true endpoint fuzzy frame} associated to \(U\) is the
pushout
\[
    \mathcal A_U
    :=
    \mathcal A
    \amalg_{\Omega_{\mathbb V}}
    \Omega_{\mathbb V}^{\top}
\]
in the category of categorical fuzzy frames.  The canonical quotient map is
denoted
\[
    j_U^\ast:\mathcal A\to\mathcal A_U.
\]
Contravariantly, this presents the generated true endpoint sublocale
\[
    j_U:U^{+}\to X.
\]
\end{definition}

\begin{definition}[Generated false endpoint fuzzy frame]
\label{def:generated-false-endpoint-frame}
The \emph{generated false endpoint fuzzy frame} associated to \(U\) is the
pushout
\[
    \mathcal A_Z
    :=
    \mathcal A
    \amalg_{\Omega_{\mathbb V}}
    \Omega_{\mathbb V}^{\bot}.
\]
The canonical quotient map is denoted
\[
    i_U^\ast:\mathcal A\to\mathcal A_Z.
\]
Contravariantly, this presents the generated false endpoint sublocale
\[
    i_U:Z_U\to X.
\]
\end{definition}

\begin{proposition}[Generated endpoint parts are fuzzy sublocales]
\label{prop:generated-endpoint-parts-are-sublocales}
The maps
\[
    j_U^\ast:\mathcal A\to\mathcal A_U,
    \qquad
    i_U^\ast:\mathcal A\to\mathcal A_Z
\]
are left-exact reflective quotients.  Hence they present fuzzy sublocales of
\(X\).
\end{proposition}

\begin{proof}
The endpoint maps
\[
    e_{\top}^{\ast}:\Omega_{\mathbb V}\to\Omega_{\mathbb V}^{\top},
    \qquad
    e_{\bot}^{\ast}:\Omega_{\mathbb V}\to\Omega_{\mathbb V}^{\bot}
\]
are left-exact reflective quotients by construction.  Base change of an
enriched finite-limit-and-scalar sketch quotient along
\[
    \chi_U^\ast:\Omega_{\mathbb V}\to\mathcal A
\]
is again an enriched finite-limit-and-scalar sketch quotient.  Therefore the
pushout maps \(j_U^\ast\) and \(i_U^\ast\) are left-exact reflective quotients
of \(\mathcal A\).
\end{proof}

\section{Generated endpoint nuclei in the posetal reflection}
\label{sec:generated-endpoint-nuclei-posetal}

Let \(A\) be a fuzzy frame in the posetal reflection, and let \(a\in A\) be the
open classified by
\[
    \Omega_\Lambda\to A.
\]

\begin{definition}[Generated true endpoint nucleus]
\label{def:generated-true-endpoint-nucleus}
The \emph{generated true endpoint nucleus} of \(a\) is
\[
    J_a^+
    :=
    \bigwedge
    \{\,j\mid j \text{ is a fuzzy nucleus on }A,\ j(a)=1\,\}.
\]
The generated true endpoint frame is
\[
    A_{J_a^+}.
\]
\end{definition}

\begin{definition}[Generated false endpoint nucleus]
\label{def:generated-false-endpoint-nucleus}
The \emph{generated false endpoint nucleus} of \(a\) is
\[
    J_a^-
    :=
    \bigwedge
    \{\,j\mid j \text{ is a fuzzy nucleus on }A,\ j(a)=j(0)\,\}.
\]
The generated false endpoint frame is
\[
    A_{J_a^-}.
\]
\end{definition}

\begin{proposition}[Universal property of the true endpoint nucleus]
\label{prop:true-endpoint-nucleus-universal}
The quotient
\[
    q_{J_a^+}:A\to A_{J_a^+}
\]
is universal among scalar-compatible fuzzy-frame quotients of \(A\) in which
\(a\) becomes top.
\end{proposition}

\begin{proof}
By construction,
\[
    J_a^+(a)=1.
\]
Hence the quotient sends \(a\) to the top element of the fixed-point frame.

Let
\[
    q:A\to B
\]
be any scalar-compatible left-exact reflective quotient sending \(a\) to top.
Let \(j_q\) be its associated fuzzy nucleus.  Since \(q(a)=1_B\), the associated
localization satisfies
\[
    j_q(a)=1.
\]
Thus \(j_q\) is one of the fuzzy nuclei appearing in the meet defining
\(J_a^+\).  Therefore
\[
    J_a^+\leq j_q.
\]
By the reversed order between nuclei and sublocales, the quotient \(q\) factors
through
\[
    q_{J_a^+}:A\to A_{J_a^+}.
\]
This proves the universal property.
\end{proof}

\begin{theorem}[The false endpoint nucleus]
\label{thm:false-endpoint-nucleus-formula}
For every \(a\in A\), the generated false endpoint nucleus is
\[
    J_a^-(x)=a\vee x.
\]
Consequently,
\[
    A_{J_a^-}\cong\uparrow a.
\]
\end{theorem}

\begin{proof}
Define
\[
    c_a(x):=a\vee x.
\]
Then \(c_a\) is inflationary and idempotent:
\[
    x\leq a\vee x,
    \qquad
    a\vee(a\vee x)=a\vee x.
\]
It preserves the terminal object and binary meets:
\[
    c_a(1)=1,
\]
and by distributivity,
\[
\begin{aligned}
    c_a(x)\wedge c_a(y)
    &=
    (a\vee x)\wedge(a\vee y)\\
    &=
    a\vee(x\wedge y)\\
    &=
    c_a(x\wedge y).
\end{aligned}
\]

It remains to check scalar compatibility.  Since the base is affine,
\[
    \lambda\leq 1,
\]
and hence
\[
    \lambda\tensor a\leq 1\tensor a=a.
\]
Therefore
\[
\begin{aligned}
    c_a(\lambda\tensor c_a(x))
    &=
    a\vee \lambda\tensor(a\vee x)\\
    &=
    a\vee(\lambda\tensor a)\vee(\lambda\tensor x)\\
    &=
    a\vee(\lambda\tensor x)\\
    &=
    c_a(\lambda\tensor x).
\end{aligned}
\]
Thus \(c_a\) is a fuzzy nucleus.

Moreover,
\[
    c_a(a)=a=c_a(0),
\]
so \(c_a\) forces \(a\) to become bottom in the quotient.

Now let \(j\) be any fuzzy nucleus satisfying
\[
    j(a)=j(0).
\]
Since \(0\leq x\), monotonicity gives
\[
    j(0)\leq j(x).
\]
Hence
\[
    a\leq j(a)=j(0)\leq j(x),
\]
and also
\[
    x\leq j(x).
\]
Therefore
\[
    a\vee x\leq j(x).
\]
Thus
\[
    c_a\leq j.
\]
So \(c_a\) is the least fuzzy nucleus forcing \(a\) to the quotient bottom:
\[
    J_a^-=c_a.
\]
The fixed points of \(c_a\) are precisely those \(x\) satisfying
\[
    a\vee x=x,
\]
equivalently
\[
    a\leq x.
\]
Thus
\[
    A_{J_a^-}\cong\uparrow a.
\]
\end{proof}

\begin{remark}[No general formula \(J_a^+(x)=a\Rightarrow x\)]
\label{rem:no-general-open-nucleus-formula}
The true endpoint nucleus
\[
    J_a^+
\]
is not generally given by the ordinary frame formula
\[
    a\Rightarrow x.
\]
That formula is valid in the ordinary/idempotent specialization, but it need
not be scalar-compatible for non-idempotent fuzzy truth-value objects.
\end{remark}

\begin{proposition}[Non-idempotent collapse phenomenon]
\label{prop:non-idempotent-true-endpoint-collapse}
Let \(A=\Lambda=[0,1]\) with the Łukasiewicz tensor
\[
    r\tensor s=\max(0,r+s-1).
\]
If \(0<a<1\), then the generated true endpoint nucleus \(J_a^+\) is the top
nucleus
\[
    x\longmapsto 1.
\]
\end{proposition}

\begin{proof}
Let \(j\) be any fuzzy nucleus satisfying
\[
    j(a)=1.
\]
Scalar compatibility gives
\[
    j(r\tensor j(a))=j(r\tensor a).
\]
Since \(j(a)=1\), this becomes
\[
    j(r)=j(r\tensor a).
\]
Taking \(r=a\) yields
\[
    j(a)=j(a\tensor a).
\]
Iterating gives
\[
    j(a)=j(a^{\tensor n})
\]
for every \(n\geq1\).

For Łukasiewicz tensor,
\[
    a^{\tensor n}=\max(0,na-(n-1)).
\]
Since \(a<1\), there exists \(n\) such that
\[
    a^{\tensor n}=0.
\]
Hence
\[
    j(0)=j(a)=1.
\]
By monotonicity, for every \(x\),
\[
    1=j(0)\leq j(x)\leq 1.
\]
Thus \(j(x)=1\) for all \(x\).  Since every fuzzy nucleus forcing \(a\) to top
is the top nucleus, their meet is the top nucleus.
\end{proof}

\section{Generated open and closed sublocales}
\label{sec:generated-open-closed-sublocales-corrected}

The preceding endpoint construction gives principal true and false endpoint
sublocales.  For the general fuzzy theory, it is useful to distinguish
principal true endpoint sublocales from arbitrary generated open sublocales.

\begin{definition}[Principal generated open sublocale]
\label{def:principal-generated-open-sublocale}
A \emph{principal generated open sublocale} of \(X\) is a sublocale obtained by
forcing a single fuzzy open
\[
    U\hookrightarrow1_{\mathcal O(X)}
\]
to be true, equivalently a quotient of the form
\[
    j_U^\ast:\mathcal O(X)\to\mathcal O(X)_U.
\]
\end{definition}

\begin{definition}[Generated open sublocale]
\label{def:generated-open-sublocale}
A \emph{generated open sublocale} of \(X\) is a join of principal generated
open sublocales.

The full subposet of generated open sublocales is denoted
\[
    \operatorname{OpenSub}(X)\subseteq\operatorname{SubLoc}(X).
\]
\end{definition}

\begin{definition}[Generated closed-complement sublocale]
\label{def:generated-closed-complement-sublocale}
A \emph{generated closed-complement sublocale} of \(X\) is a sublocale obtained
by forcing a single fuzzy open to be false.

The full subposet of generated closed-complement sublocales is denoted
\[
    \operatorname{ClosedSub}(X)\subseteq\operatorname{SubLoc}(X).
\]
\end{definition}

\begin{proposition}[Order of principal closed complements]
\label{prop:order-principal-closed-complements}
In the posetal reflection, for \(a,b\in A\),
\[
    a\leq b
    \quad\Longleftrightarrow\quad
    Z_b\leq Z_a.
\]
Consequently, principal closed-complement sublocales are ordered oppositely to
their generating opens.
\end{proposition}

\begin{proof}
The closed-complement sublocale \(Z_a\) is presented by the fuzzy nucleus
\[
    J_a^-(x)=a\vee x.
\]
By Proposition~\ref{prop:sublocale-order-nuclei},
\[
    Z_b\leq Z_a
    \quad\Longleftrightarrow\quad
    J_a^-\leq J_b^-.
\]
This means
\[
    a\vee x\leq b\vee x
\]
for all \(x\).  If \(a\leq b\), this is immediate.  Conversely, taking \(x=0\)
gives
\[
    a\leq b.
\]
\end{proof}

\section{Adjoint calculus and attachment}
\label{sec:adjoint-calculus-and-attachment-corrected}

Let
\[
    j_U^\ast:\mathcal A\to\mathcal A_U
\]
be the generated true endpoint quotient.  Since it is a reflective quotient,
there is a fully faithful right adjoint
\[
    (j_U)_\ast:\mathcal A_U\to\mathcal A,
\]
with
\[
    j_U^\ast\dashv(j_U)_\ast.
\]

Similarly, for the generated false endpoint quotient
\[
    i_U^\ast:\mathcal A\to\mathcal A_Z,
\]
there is a fully faithful right adjoint
\[
    (i_U)_\ast:\mathcal A_Z\to\mathcal A,
\]
with
\[
    i_U^\ast\dashv(i_U)_\ast.
\]

\begin{definition}[Attachment functor]
\label{def:fuzzy-attachment-functor}
The \emph{attachment functor} of the generated endpoint pair is
\[
    \Theta_U:\mathcal A_U\to\mathcal A_Z
\]
defined by
\[
    \Theta_U:=i_U^\ast (j_U)_\ast.
\]
\end{definition}

\begin{proposition}[Attachment is left exact]
\label{prop:attachment-left-exact}
The attachment functor
\[
    \Theta_U:\mathcal A_U\to\mathcal A_Z
\]
preserves the specified finite weighted limits.
\end{proposition}

\begin{proof}
The right adjoint
\[
    (j_U)_\ast:\mathcal A_U\to\mathcal A
\]
preserves finite weighted limits.  The quotient
\[
    i_U^\ast:\mathcal A\to\mathcal A_Z
\]
is left exact by construction.  Therefore the composite
\[
    \Theta_U=i_U^\ast(j_U)_\ast
\]
preserves the specified finite weighted limits.
\end{proof}

\begin{proposition}[Attachment in the posetal reflection]
\label{prop:attachment-posetal-general}
In the posetal reflection, the attachment functor has the form
\[
    \Theta_a:A_{J_a^+}\to A_{J_a^-},
    \qquad
    \Theta_a(h)=a\vee h,
\]
where \(h\in A_{J_a^+}\subseteq A\).
\end{proposition}

\begin{proof}
The right adjoint
\[
    (j_a)_\ast:A_{J_a^+}\hookrightarrow A
\]
is the inclusion of fixed points.  The quotient
\[
    i_a^\ast:A\to A_{J_a^-}
\]
is given by the false endpoint nucleus
\[
    J_a^-(x)=a\vee x.
\]
Therefore
\[
    \Theta_a(h)=i_a^\ast(h)=J_a^-(h)=a\vee h.
\]
\end{proof}

\section{Artin gluing of categorical fuzzy frames}
\label{sec:artin-gluing-categorical-fuzzy-frames-corrected}

The attachment functor is left exact, but it need not be a fuzzy-frame morphism
in the strong sense of preserving all small weighted colimits.  Artin gluing is
therefore formulated for left-exact enriched functors.

\begin{definition}[Artin gluing]
\label{def:enriched-artin-gluing}
Let
\[
    \Theta:\mathcal H\to\mathcal N
\]
be a left-exact \(\mathbb V\)-functor between categorical fuzzy frames.  The
\emph{Artin gluing} of \(\Theta\) is the enriched comma fuzzy frame
\[
    \operatorname{Gl}(\Theta)
    :=
    (\id_{\mathcal N}\downarrow\Theta).
\]
Its objects are triples
\[
    (n,h,\alpha),
\]
where
\[
    n\in\mathcal N,
    \qquad
    h\in\mathcal H,
    \qquad
    \alpha:n\to\Theta(h)
\]
is a morphism in \(\mathcal N\).
\end{definition}

The projection functors are denoted
\[
    \pi_{\mathcal N}:\operatorname{Gl}(\Theta)\to\mathcal N,
    \qquad
    \pi_{\mathcal H}:\operatorname{Gl}(\Theta)\to\mathcal H.
\]

\begin{proposition}[Gluing is a categorical fuzzy frame]
\label{prop:gluing-is-fuzzy-frame}
If \(\mathcal H\) and \(\mathcal N\) are categorical fuzzy frames and
\[
    \Theta:\mathcal H\to\mathcal N
\]
preserves the specified finite weighted limits, then
\[
    \operatorname{Gl}(\Theta)
\]
is a categorical fuzzy frame.
\end{proposition}

\begin{proof}
The enriched comma object exists because the relevant enriched limits exist in
the ambient universe.

Finite weighted limits are computed componentwise.  For a finite weighted
diagram
\[
    k\longmapsto (n_k,h_k,\alpha_k),
\]
the limit has first two components
\[
    \lim n_k,
    \qquad
    \lim h_k.
\]
The structure map
\[
    \lim n_k\to \Theta(\lim h_k)
\]
is obtained from the cone
\[
    \lim n_k\to n_k\xrightarrow{\alpha_k}\Theta(h_k)
\]
together with the left-exactness isomorphism
\[
    \Theta(\lim h_k)\cong\lim\Theta(h_k).
\]

Small weighted colimits are also computed componentwise.  For a weighted
diagram
\[
    k\longmapsto(n_k,h_k,\alpha_k)
\]
with weight \(W\), the colimit has first two components
\[
    W*n_k,
    \qquad
    W*h_k.
\]
The structure map
\[
    W*n_k\to\Theta(W*h_k)
\]
is induced by the weighted-colimit universal property from the composites
\[
    n_k
    \xrightarrow{\alpha_k}
    \Theta(h_k)
    \longrightarrow
    \Theta(W*h_k),
\]
where the final arrow is obtained by applying \(\Theta\) to the colimit cocone
\[
    h_k\to W*h_k.
\]
The universal property of the comma object then gives the universal property
of the weighted colimit in \(\operatorname{Gl}(\Theta)\).

Thus \(\operatorname{Gl}(\Theta)\) has the required small weighted colimits and
specified finite weighted limits, and these are compatible with the enriched
structure.  Hence it is a categorical fuzzy frame.
\end{proof}

\begin{proposition}[Posetal Artin gluing]
\label{prop:posetal-artin-gluing}
In the posetal reflection, a left-exact map
\[
    \theta:H\to N
\]
between fuzzy frames has Artin gluing
\[
    \operatorname{Gl}(\theta)
    =
    \{(n,h)\in N\times H\mid n\leq\theta(h)\}.
\]
Finite meets and arbitrary joins are computed componentwise.
\end{proposition}

\begin{proof}
If
\[
    n\leq\theta(h),
    \qquad
    n'\leq\theta(h'),
\]
then left exactness gives
\[
    n\wedge n'
    \leq
    \theta(h)\wedge\theta(h')
    =
    \theta(h\wedge h').
\]
Thus componentwise finite meets remain in the gluing frame.

For joins, suppose
\[
    n_i\leq\theta(h_i)
\]
for all \(i\).  The cocone maps
\[
    h_i\leq\bigvee_i h_i
\]
give
\[
    \theta(h_i)\leq\theta\left(\bigvee_i h_i\right).
\]
Therefore
\[
    n_i\leq\theta\left(\bigvee_i h_i\right)
\]
for all \(i\), and hence
\[
    \bigvee_i n_i
    \leq
    \theta\left(\bigvee_i h_i\right).
\]
Thus componentwise joins remain in the gluing frame.  The universal properties
are inherited from the product order.
\end{proof}

\section{The canonical endpoint-gluing comparison}
\label{sec:canonical-endpoint-gluing-comparison}

Let
\[
    U\hookrightarrow1_{\mathcal A}
\]
be a fuzzy open, with generated endpoint quotients
\[
    j_U^\ast:\mathcal A\to\mathcal A_U,
    \qquad
    i_U^\ast:\mathcal A\to\mathcal A_Z.
\]
Let
\[
    \Theta_U=i_U^\ast(j_U)_\ast:
    \mathcal A_U\to\mathcal A_Z
\]
be the attachment functor.

\begin{construction}[The endpoint-gluing comparison]
\label{con:endpoint-gluing-comparison}
There is a canonical fuzzy-frame morphism
\[
    \Phi_U:\mathcal A\to\operatorname{Gl}(\Theta_U)
\]
defined on objects by
\[
    \Phi_U(X)
    =
    \left(
        i_U^\ast X,\,
        j_U^\ast X,\,
        i_U^\ast X\to i_U^\ast(j_U)_\ast j_U^\ast X
    \right).
\]
The structure map is obtained by applying \(i_U^\ast\) to the unit of the
adjunction
\[
    j_U^\ast\dashv(j_U)_\ast:
    \qquad
    X\to(j_U)_\ast j_U^\ast X.
\]
\end{construction}

\begin{proposition}[The endpoint-gluing comparison is canonical]
\label{prop:endpoint-gluing-comparison-canonical}
The construction above defines a canonical morphism of categorical fuzzy
frames
\[
    \Phi_U:\mathcal A\to\operatorname{Gl}(\Theta_U).
\]
\end{proposition}

\begin{proof}
Both components
\[
    i_U^\ast:\mathcal A\to\mathcal A_Z,
    \qquad
    j_U^\ast:\mathcal A\to\mathcal A_U
\]
are fuzzy-frame morphisms.  The structure map is natural in \(X\) because it is
obtained from the unit of an enriched adjunction.  Compatibility with the
specified finite weighted limits follows from left exactness of \(i_U^\ast\) and
\(j_U^\ast\).  Compatibility with small weighted colimits follows from the
fact that the two endpoint quotient maps preserve such colimits and gluing
colimits are computed componentwise.
\end{proof}

\begin{proposition}[Posetal form of the endpoint-gluing comparison]
\label{prop:posetal-endpoint-gluing-comparison}
In the posetal reflection, the endpoint-gluing comparison is
\[
    \Phi_a:A\to\operatorname{Gl}(\Theta_a),
    \qquad
    \Phi_a(x)=\bigl(J_a^-(x),J_a^+(x)\bigr).
\]
The gluing condition is
\[
    J_a^-(x)\leq J_a^-(J_a^+(x)).
\]
\end{proposition}

\begin{proof}
The first component is the false endpoint quotient
\[
    x\mapsto J_a^-(x),
\]
and the second component is the true endpoint quotient
\[
    x\mapsto J_a^+(x).
\]
The attachment map is
\[
    \Theta_a(h)=J_a^-(h).
\]
Since
\[
    x\leq J_a^+(x)
\]
and \(J_a^-\) is monotone,
\[
    J_a^-(x)\leq J_a^-(J_a^+(x)).
\]
Thus the displayed pair lies in the gluing frame.
\end{proof}

\begin{remark}[No general reconstruction theorem]
\label{rem:no-general-reconstruction-theorem}
The comparison
\[
    \Phi_U:\mathcal A\to\operatorname{Gl}(\Theta_U)
\]
is not generally an equivalence.  For non-idempotent bases, the true endpoint
nucleus may collapse, as in Proposition~\ref{prop:non-idempotent-true-endpoint-collapse}.
Thus the ordinary open--closed reconstruction theorem cannot hold in the full
scalar-compatible fuzzy setting.
\end{remark}

\section{Interior and closure of fuzzy sublocales}
\label{sec:fuzzy-interior-closure-corrected}

Let \(X\) be a fuzzy locale.

\begin{definition}[Interior]
\label{def:fuzzy-sublocale-interior}
Let
\[
    S\in\operatorname{SubLoc}(X).
\]
The \emph{interior} of \(S\) is the largest generated open sublocale contained
in \(S\):
\[
    \operatorname{int}(S)
    :=
    \bigvee
    \{\,U\in\operatorname{OpenSub}(X)\mid U\leq S\,\}.
\]
\end{definition}

\begin{definition}[Closure]
\label{def:fuzzy-sublocale-closure}
Let
\[
    S\in\operatorname{SubLoc}(X).
\]
The \emph{closure} of \(S\) is the smallest generated closed-complement
sublocale containing \(S\):
\[
    \operatorname{cl}(S)
    :=
    \bigwedge
    \{\,Z\in\operatorname{ClosedSub}(X)\mid S\leq Z\,\}.
\]
\end{definition}

Joins and meets of fuzzy sublocales are computed in the lattice of left-exact
reflective quotients by combining the corresponding enriched
finite-limit-and-scalar quotient data.

\begin{proposition}[Interior adjunction]
\label{prop:fuzzy-interior-adjunction}
The inclusion
\[
    \iota_{\mathrm{open}}:
    \operatorname{OpenSub}(X)
    \hookrightarrow
    \operatorname{SubLoc}(X)
\]
has right adjoint
\[
    \operatorname{int}.
\]
Thus
\[
    \iota_{\mathrm{open}}\dashv\operatorname{int}.
\]
\end{proposition}

\begin{proof}
By definition, \(\operatorname{int}(S)\) is the join of all generated open
sublocales below \(S\).  Hence, for a generated open sublocale \(U\),
\[
    U\leq S
    \quad\Longleftrightarrow\quad
    U\leq\operatorname{int}(S).
\]
This is precisely the adjunction.
\end{proof}

\begin{proposition}[Closure adjunction]
\label{prop:fuzzy-closure-adjunction}
The inclusion
\[
    \iota_{\mathrm{closed}}:
    \operatorname{ClosedSub}(X)
    \hookrightarrow
    \operatorname{SubLoc}(X)
\]
has left adjoint
\[
    \operatorname{cl}.
\]
Thus
\[
    \operatorname{cl}\dashv\iota_{\mathrm{closed}}.
\]
\end{proposition}

\begin{proof}
By definition, \(\operatorname{cl}(S)\) is the meet of all generated
closed-complement sublocales above \(S\).  Hence, for a generated closed
sublocale \(Z\),
\[
    \operatorname{cl}(S)\leq Z
    \quad\Longleftrightarrow\quad
    S\leq Z.
\]
This is precisely the adjunction.
\end{proof}

\begin{corollary}[Interior comonad and closure monad]
\label{cor:fuzzy-interior-comonad-closure-monad}
The composite
\[
    \mathbb I:=\iota_{\mathrm{open}}\operatorname{int}
\]
is an idempotent comonad on \(\operatorname{SubLoc}(X)\), and the composite
\[
    \mathbb C:=\iota_{\mathrm{closed}}\operatorname{cl}
\]
is an idempotent monad on \(\operatorname{SubLoc}(X)\).  Moreover,
\[
    \mathbb I(S)\leq S\leq\mathbb C(S)
\]
for every fuzzy sublocale \(S\).
\end{corollary}

\begin{proof}
The first statement follows from
\[
    \iota_{\mathrm{open}}\dashv\operatorname{int}
\]
and full faithfulness of the inclusion of generated open sublocales.  The
second follows from
\[
    \operatorname{cl}\dashv\iota_{\mathrm{closed}}
\]
and full faithfulness of the inclusion of generated closed-complement
sublocales.  The displayed inequalities are the counit and unit inequalities of
these two adjunctions.
\end{proof}

\begin{theorem}[Closure formula in the posetal reflection]
\label{thm:posetal-closure-formula}
Let \(X_\nu\) be a fuzzy sublocale presented by a fuzzy nucleus
\[
    \nu:A\to A.
\]
Then
\[
    \operatorname{cl}(X_\nu)=Z_{\nu(0)}.
\]
Equivalently, the closure is presented by the fuzzy nucleus
\[
    x\longmapsto\nu(0)\vee x.
\]
\end{theorem}

\begin{proof}
A generated closed-complement sublocale \(Z_a\) contains \(X_\nu\) precisely
when
\[
    X_\nu\leq Z_a.
\]
By order reversal, this is equivalent to
\[
    J_a^-\leq\nu.
\]
Using the formula
\[
    J_a^-(x)=a\vee x,
\]
this means
\[
    a\vee x\leq\nu(x)
\]
for every \(x\in A\).  Taking \(x=0\) gives
\[
    a\leq\nu(0).
\]

Conversely, if \(a\leq\nu(0)\), then for every \(x\),
\[
    a\vee x
    \leq
    \nu(0)\vee x
    \leq
    \nu(x),
\]
using monotonicity and inflationarity of \(\nu\).  Thus
\[
    J_a^-\leq\nu.
\]
Therefore the smallest generated closed-complement sublocale containing
\(X_\nu\) is
\[
    Z_{\nu(0)}.
\]
\end{proof}

\section{The ordinary and idempotent specialization}
\label{sec:ordinary-idempotent-specialization}

We now recover the classical formulas under the idempotent/cartesian
specialization.  In this section the scalar tensor acts as meet:
\[
    \lambda\tensor x=\lambda\wedge x.
\]
This includes the ordinary crisp case \(\Lambda=2\), and also the ordinary
underlying-frame reflection where scalar compatibility is no additional
restriction.

\begin{theorem}[Idempotent true and false endpoint nuclei]
\label{thm:idempotent-endpoint-nuclei}
In the idempotent specialization, the generated endpoint nuclei of an open
\(a\in A\) are
\[
    J_a^+(x)=a\Rightarrow x,
    \qquad
    J_a^-(x)=a\vee x.
\]
\end{theorem}

\begin{proof}
The formula for \(J_a^-\) is Theorem~\ref{thm:false-endpoint-nucleus-formula}.

Define
\[
    o_a(x):=a\Rightarrow x.
\]
This is the usual open nucleus of the underlying frame.  It is inflationary
because
\[
    a\wedge x\leq x.
\]
It preserves finite meets because it is right adjoint to
\[
    a\wedge-.
\]
It is idempotent since
\[
    a\Rightarrow(a\Rightarrow x)
    =
    (a\wedge a)\Rightarrow x
    =
    a\Rightarrow x.
\]

We check scalar compatibility.  Since the scalar tensor is meet,
\[
\begin{aligned}
    o_a(\lambda\wedge o_a(x))
    &=
    a\Rightarrow(\lambda\wedge(a\Rightarrow x))\\
    &=
    (a\Rightarrow\lambda)\wedge(a\Rightarrow(a\Rightarrow x))\\
    &=
    (a\Rightarrow\lambda)\wedge(a\Rightarrow x)\\
    &=
    a\Rightarrow(\lambda\wedge x)\\
    &=
    o_a(\lambda\wedge x).
\end{aligned}
\]
Thus \(o_a\) is a fuzzy nucleus in the idempotent specialization.

Moreover,
\[
    o_a(a)=a\Rightarrow a=1.
\]
Let \(j\) be any fuzzy nucleus with \(j(a)=1\).  Since
\[
    a\wedge(a\Rightarrow x)\leq x,
\]
monotonicity gives
\[
    j(a\wedge(a\Rightarrow x))\leq j(x).
\]
Since \(j\) preserves binary meets,
\[
    j(a\wedge(a\Rightarrow x))
    =
    j(a)\wedge j(a\Rightarrow x)
    =
    1\wedge j(a\Rightarrow x)
    =
    j(a\Rightarrow x).
\]
By inflationarity,
\[
    a\Rightarrow x\leq j(a\Rightarrow x).
\]
Therefore
\[
    a\Rightarrow x\leq j(x).
\]
Thus
\[
    o_a\leq j
\]
for every fuzzy nucleus \(j\) forcing \(a\) to top.  Hence
\[
    J_a^+=o_a.
\]
\end{proof}

\begin{proposition}[Generated open and closed frames in the idempotent case]
\label{prop:idempotent-open-closed-frames}
In the idempotent specialization,
\[
    A_{J_a^+}\cong\downarrow a,
    \qquad
    A_{J_a^-}\cong\uparrow a.
\]
The quotient maps are
\[
    j_a^\ast:A\to\downarrow a,
    \qquad
    j_a^\ast(x)=a\wedge x,
\]
and
\[
    i_a^\ast:A\to\uparrow a,
    \qquad
    i_a^\ast(x)=a\vee x.
\]
\end{proposition}

\begin{proof}
The fixed points of
\[
    J_a^-(x)=a\vee x
\]
are precisely the elements \(x\) satisfying \(a\leq x\), so
\[
    A_{J_a^-}\cong\uparrow a.
\]

For the open part, the fixed points of
\[
    J_a^+(x)=a\Rightarrow x
\]
are isomorphic to \(\downarrow a\).  The maps are
\[
    A_{J_a^+}\to\downarrow a,
    \qquad
    h\longmapsto a\wedge h,
\]
and
\[
    \downarrow a\to A_{J_a^+},
    \qquad
    u\longmapsto a\Rightarrow u.
\]
If \(u\leq a\), then
\[
    a\wedge(a\Rightarrow u)=u.
\]
If \(h=a\Rightarrow h\), then
\[
    a\Rightarrow(a\wedge h)
    =
    (a\Rightarrow a)\wedge(a\Rightarrow h)
    =
    h.
\]
Thus the two maps are inverse.

Under this identification, the quotient
\[
    q_{J_a^+}:A\to A_{J_a^+}
\]
corresponds to
\[
    x\longmapsto a\wedge x.
\]
The false quotient is already
\[
    x\longmapsto a\vee x.
\]
\end{proof}

\begin{proposition}[Open immersion adjoints in the idempotent case]
\label{prop:idempotent-open-adjoints}
The map
\[
    j_a^\ast:A\to\downarrow a,
    \qquad
    j_a^\ast(x)=a\wedge x,
\]
has both a left and a right adjoint:
\[
    (j_a)_!\dashv j_a^\ast\dashv(j_a)_\ast.
\]
They are
\[
    (j_a)_!(u)=u,
    \qquad
    (j_a)_\ast(u)=a\Rightarrow u.
\]
\end{proposition}

\begin{proof}
For \(u\leq a\) and \(x\in A\),
\[
    u\leq x
    \quad\Longleftrightarrow\quad
    u\leq a\wedge x.
\]
Thus
\[
    (j_a)_!\dashv j_a^\ast.
\]
Also,
\[
    a\wedge x\leq u
    \quad\Longleftrightarrow\quad
    x\leq a\Rightarrow u.
\]
Thus
\[
    j_a^\ast\dashv(j_a)_\ast.
\]
\end{proof}

\begin{proposition}[Closed immersion adjoint in the idempotent case]
\label{prop:idempotent-closed-adjoint}
The generated closed-complement inverse image
\[
    i_a^\ast:A\to\uparrow a,
    \qquad
    i_a^\ast(x)=a\vee x,
\]
has right adjoint given by inclusion:
\[
    (i_a)_\ast:\uparrow a\hookrightarrow A.
\]
\end{proposition}

\begin{proof}
For \(x\in A\) and \(n\in\uparrow a\),
\[
    a\vee x\leq n
    \quad\Longleftrightarrow\quad
    x\leq n,
\]
because \(a\leq n\).  This is the adjunction
\[
    i_a^\ast\dashv(i_a)_\ast.
\]
\end{proof}

\begin{proposition}[Attachment formula in the idempotent case]
\label{prop:idempotent-attachment-formula}
Under the identifications
\[
    A_{J_a^+}\cong\downarrow a,
    \qquad
    A_{J_a^-}\cong\uparrow a,
\]
the attachment map is
\[
    \theta_a:\downarrow a\to\uparrow a,
    \qquad
    \theta_a(u)=a\vee(a\Rightarrow u).
\]
\end{proposition}

\begin{proof}
In general,
\[
    \Theta_a(h)=a\vee h
\]
for \(h\in A_{J_a^+}\subseteq A\).  Under the isomorphism
\[
    \downarrow a\cong A_{J_a^+},
\]
an element \(u\leq a\) corresponds to
\[
    a\Rightarrow u.
\]
Therefore
\[
    \theta_a(u)=a\vee(a\Rightarrow u).
\]
\end{proof}

\section{Open--closed complementarity in the idempotent case}
\label{sec:idempotent-open-closed-complementarity}

\begin{theorem}[Generated open and closed complement are complementary]
\label{thm:idempotent-open-closed-complementary}
In the idempotent specialization, the generated open sublocale \(U_a\) and its
generated closed complement \(Z_a\) satisfy
\[
    U_a\vee Z_a=X,
    \qquad
    U_a\wedge Z_a=0.
\]
\end{theorem}

\begin{proof}
The open and closed sublocales are presented by the nuclei
\[
    O_a(x)=a\Rightarrow x,
    \qquad
    C_a(x)=a\vee x.
\]

The join of sublocales corresponds to the pointwise meet of nuclei.  For every
\(x\),
\[
\begin{aligned}
    O_a(x)\wedge C_a(x)
    &=
    (a\Rightarrow x)\wedge(a\vee x)\\
    &=
    \bigl(a\wedge(a\Rightarrow x)\bigr)
    \vee
    \bigl(x\wedge(a\Rightarrow x)\bigr)\\
    &=
    x.
\end{aligned}
\]
Thus
\[
    O_a\wedge C_a=\id_A,
\]
and therefore
\[
    U_a\vee Z_a=X.
\]

For the meet, let \(\nu\) be any nucleus such that
\[
    O_a\leq\nu,
    \qquad
    C_a\leq\nu.
\]
Since
\[
    C_a(0)=a,
\]
we have
\[
    a\leq\nu(0).
\]
Therefore
\[
    O_a(\nu(0))=a\Rightarrow\nu(0)=1.
\]
Using \(O_a\leq\nu\) and idempotence of \(\nu\), we get
\[
    1=O_a(\nu(0))\leq\nu(\nu(0))=\nu(0).
\]
Thus
\[
    \nu(0)=1.
\]
A nucleus satisfying \(\nu(0)=1\) is the top nucleus.  Hence the meet of
\(U_a\) and \(Z_a\) is the bottom sublocale:
\[
    U_a\wedge Z_a=0.
\]
\end{proof}

\section{Open--closed Artin gluing reconstruction in the idempotent case}
\label{sec:idempotent-open-closed-artin-gluing}

Let \(A\) be a frame in the idempotent specialization, and let \(a\in A\).  The
generated open and closed parts are
\[
    \downarrow a,
    \qquad
    \uparrow a,
\]
and the attachment map is
\[
    \theta_a(u)=a\vee(a\Rightarrow u).
\]

\begin{theorem}[Open--closed Artin gluing reconstruction]
\label{thm:idempotent-open-closed-gluing-reconstruction}
There is a canonical isomorphism of frames
\[
    A
    \cong
    \operatorname{Gl}(\theta_a)
    =
    \{(n,u)\in\uparrow a\times\downarrow a
      \mid n\leq a\vee(a\Rightarrow u)\}.
\]
The forward map is
\[
    \Phi(x)=(a\vee x,\ a\wedge x),
\]
and the inverse map is
\[
    \Psi(n,u)=n\wedge(a\Rightarrow u).
\]
\end{theorem}

\begin{proof}
First we show that \(\Phi(x)\) lies in the gluing frame.  We need
\[
    a\vee x
    \leq
    a\vee(a\Rightarrow(a\wedge x)).
\]
It is enough to show
\[
    x\leq a\Rightarrow(a\wedge x),
\]
which is equivalent to
\[
    a\wedge x\leq a\wedge x.
\]
Thus \(\Phi\) is well defined.

Now compute
\[
\begin{aligned}
    \Psi\Phi(x)
    &=
    (a\vee x)\wedge(a\Rightarrow(a\wedge x)).
\end{aligned}
\]
The inequality
\[
    x\leq\Psi\Phi(x)
\]
follows from the preceding paragraph.  Conversely,
\[
\begin{aligned}
    (a\vee x)\wedge(a\Rightarrow(a\wedge x))
    &=
    \bigl(a\wedge(a\Rightarrow(a\wedge x))\bigr)
    \vee
    \bigl(x\wedge(a\Rightarrow(a\wedge x))\bigr)\\
    &\leq
    (a\wedge x)\vee x\\
    &=
    x.
\end{aligned}
\]
Therefore
\[
    \Psi\Phi=\id_A.
\]

Conversely, let
\[
    (n,u)\in\operatorname{Gl}(\theta_a).
\]
Thus
\[
    a\leq n,
    \qquad
    u\leq a,
    \qquad
    n\leq a\vee(a\Rightarrow u).
\]
Then
\[
\begin{aligned}
    a\vee\Psi(n,u)
    &=
    a\vee(n\wedge(a\Rightarrow u))\\
    &=
    (a\vee n)\wedge(a\vee(a\Rightarrow u))\\
    &=
    n\wedge(a\vee(a\Rightarrow u))\\
    &=
    n.
\end{aligned}
\]
Also,
\[
\begin{aligned}
    a\wedge\Psi(n,u)
    &=
    a\wedge n\wedge(a\Rightarrow u)\\
    &=
    a\wedge(a\Rightarrow u)\\
    &=
    u,
\end{aligned}
\]
using \(a\leq n\) and \(u\leq a\).  Hence
\[
    \Phi\Psi=\id_{\operatorname{Gl}(\theta_a)}.
\]
Thus \(\Phi\) and \(\Psi\) are inverse isomorphisms.
\end{proof}

\section{Interior formula in the idempotent case}
\label{sec:idempotent-interior-formula}

Let \(X_\nu\) be a sublocale presented by a nucleus
\[
    \nu:A\to A.
\]

\begin{theorem}[Interior formula in the idempotent case]
\label{thm:idempotent-interior-formula}
In the idempotent specialization,
\[
    \operatorname{int}(X_\nu)
    =
    U_{\bigwedge_{x\in A}(\nu(x)\Rightarrow x)}.
\]
\end{theorem}

\begin{proof}
An open sublocale \(U_a\) is contained in \(X_\nu\) precisely when
\[
    U_a\leq X_\nu.
\]
By order reversal, this is equivalent to
\[
    \nu\leq O_a,
\]
where
\[
    O_a(x)=a\Rightarrow x.
\]
Thus, for every \(x\in A\),
\[
    \nu(x)\leq a\Rightarrow x.
\]
By the Heyting adjunction, this is equivalent to
\[
    a\wedge\nu(x)\leq x,
\]
or
\[
    a\leq\nu(x)\Rightarrow x.
\]
Therefore \(U_a\leq X_\nu\) holds precisely when
\[
    a\leq
    \bigwedge_{x\in A}(\nu(x)\Rightarrow x).
\]
Since generated open sublocales are ordered by their generating opens in the
idempotent case, the largest such generated open is
\[
    U_{\bigwedge_{x\in A}(\nu(x)\Rightarrow x)}.
\]
\end{proof}

\section{Predicate-frame and fuzzy-domain instances}
\label{sec:predicate-frame-fuzzy-domain-instances-corrected}

The general constructions of this chapter apply to every categorical fuzzy
frame constructed earlier.  In particular, they apply to fuzzy predicate
frames, localized fuzzy predicate frames, and fuzzy Scott frames of fuzzy
domains.

\begin{proposition}[Predicate-frame and Scott-frame instances]
\label{prop:predicate-frame-scott-frame-instances}
Let \(\mathcal A\) be any categorical fuzzy frame obtained from a fuzzy
predicate construction, a localized fuzzy predicate construction, or a fuzzy
Scott-open construction.  Then \(\mathcal A\) admits:

\begin{enumerate}
    \item fuzzy sublocales;
    \item generated true endpoint sublocales;
    \item generated false endpoint sublocales;
    \item attachment functors;
    \item Artin gluing frames;
    \item canonical endpoint-gluing comparisons;
    \item interior and closure of fuzzy sublocales.
\end{enumerate}
\end{proposition}

\begin{proof}
Each item is defined for arbitrary categorical fuzzy frames.  Since the fuzzy
predicate and fuzzy Scott constructions produce categorical fuzzy frames, all
the preceding constructions apply.
\end{proof}

\begin{remark}[Sublocales are localic, not necessarily spatial]
\label{rem:sublocales-localic-not-spatial}
When \(\mathcal A\) is a predicate frame associated to a fuzzy category,
sublocales of the locale presented by \(\mathcal A\) should not automatically be
identified with subobjects or subcategories of the original fuzzy category.
Such an identification requires an additional representation theorem.
\end{remark}

\section{Crisp separated specialization}
\label{sec:crisp-separated-open-closed-gluing}

The crisp separated case is obtained by taking
\[
    \Lambda=2
\]
and restricting to separated fuzzy frames and fuzzy locales.  Categorical fuzzy
frames become ordinary frames, fuzzy locales become ordinary locales, and fuzzy
opens become ordinary opens.

For an ordinary frame \(A\) and an open \(a\in A\), the preceding idempotent
formulas become:
\[
    O_a(x)=a\Rightarrow x,
    \qquad
    C_a(x)=a\vee x.
\]
The frames of the generated open and closed-complement sublocales are
\[
    \mathcal O(U_a)\cong\downarrow a,
    \qquad
    \mathcal O(Z_a)\cong\uparrow a.
\]
The inverse-image frame maps are
\[
    j_a^\ast(x)=a\wedge x,
    \qquad
    i_a^\ast(x)=a\vee x.
\]
The open immersion has adjoint triple
\[
    (j_a)_!\dashv j_a^\ast\dashv(j_a)_\ast,
\]
where
\[
    (j_a)_!(u)=u,
    \qquad
    (j_a)_\ast(u)=a\Rightarrow u.
\]
The closed-complement immersion has
\[
    i_a^\ast\dashv(i_a)_\ast,
\]
where \((i_a)_\ast\) is inclusion.

The attachment map is
\[
    \theta_a(u)=a\vee(a\Rightarrow u),
\]
and the open--closed Artin gluing reconstruction is
\[
    A\cong
    \{(n,u)\in\uparrow a\times\downarrow a
      \mid
      n\leq a\vee(a\Rightarrow u)\}.
\]
The forward and inverse maps are
\[
    x\longmapsto(a\vee x,\ a\wedge x),
\]
and
\[
    (n,u)\longmapsto n\wedge(a\Rightarrow u).
\]

Thus the ordinary localic open--closed decomposition and Artin gluing theorem
are recovered exactly as the crisp separated specialization of the endpoint
construction.

\section{Summary}
\label{sec:summary-fuzzy-sublocales-open-closed-gluing}

This chapter developed fuzzy sublocales and endpoint gluing at the level of
categorical fuzzy frames.

A fuzzy sublocale of \(X\), with
\[
    \mathcal O(X)=\mathcal A,
\]
is presented by a left-exact reflective quotient
\[
    q\dashv r:
    \mathcal A
    \rightleftarrows
    \mathcal B.
\]
In the posetal reflection, such quotients are exactly scalar-compatible fuzzy
nuclei
\[
    j:A\to A.
\]

A fuzzy open
\[
    U\hookrightarrow1_{\mathcal A}
\]
is classified by
\[
    \chi_U^\ast:\Omega_{\mathbb V}\to\mathcal A.
\]
The generated true and false endpoint frames are obtained by endpoint base
change:
\[
    \mathcal A_U
    =
    \mathcal A\amalg_{\Omega_{\mathbb V}}\Omega_{\mathbb V}^{\top},
    \qquad
    \mathcal A_Z
    =
    \mathcal A\amalg_{\Omega_{\mathbb V}}\Omega_{\mathbb V}^{\bot}.
\]

In the posetal scalar-compatible reflection, the true endpoint is presented by
the generated fuzzy nucleus
\[
    J_a^+
    =
    \bigwedge
    \{\,j\mid j(a)=1\,\},
\]
while the false endpoint has the explicit formula
\[
    J_a^-(x)=a\vee x.
\]
The true endpoint is not generally given by the ordinary formula
\[
    a\Rightarrow x.
\]
That formula is valid only in the idempotent, ordinary, or crisp
specialization.

The attachment functor is
\[
    \Theta_U=i_U^\ast(j_U)_\ast:
    \mathcal A_U\to\mathcal A_Z.
\]
In the posetal reflection it is
\[
    \Theta_a(h)=a\vee h
\]
for \(h\in A_{J_a^+}\).  Artin gluing is defined for any left-exact enriched
functor
\[
    \Theta:\mathcal H\to\mathcal N
\]
by
\[
    \operatorname{Gl}(\Theta)
    =
    (\id_{\mathcal N}\downarrow\Theta).
\]
For every fuzzy open there is a canonical endpoint-gluing comparison
\[
    \mathcal A\to\operatorname{Gl}(\Theta_U).
\]
This comparison is not an equivalence in general non-idempotent fuzzy bases.

Closure is fully general.  If \(X_\nu\) is presented by a fuzzy nucleus
\[
    \nu:A\to A,
\]
then
\[
    \operatorname{cl}(X_\nu)=Z_{\nu(0)}.
\]
Interior is defined abstractly as the largest generated open sublocale below a
given sublocale.

In the idempotent specialization, the classical formulas are recovered:
\[
    J_a^+(x)=a\Rightarrow x,
    \qquad
    J_a^-(x)=a\vee x,
\]
\[
    A_{J_a^+}\cong\downarrow a,
    \qquad
    A_{J_a^-}\cong\uparrow a,
\]
and
\[
    \theta_a(u)=a\vee(a\Rightarrow u).
\]
The generated open and closed-complement sublocales are complementary, and the
ordinary open--closed Artin gluing reconstruction holds:
\[
    A\cong
    \{(n,u)\in\uparrow a\times\downarrow a
      \mid
      n\leq a\vee(a\Rightarrow u)\}.
\]
The inverse map is
\[
    (n,u)\longmapsto n\wedge(a\Rightarrow u).
\]

Thus the crisp localic chapter is recovered as the idempotent reflection of
the more general endpoint-quotient calculus for categorical fuzzy frames.